
\section{定界余链的规范等价}
\label{sec:4-3}

本节在定界余链集合 $\widehat{\mathcal M}(C)$ 上定义等价关系
$\sim$，并证明商
\[
 \mathcal M(C)=\widehat{\mathcal M}(C)/\mathord\sim
\]
是滤过 $A_\infty$ 代数 $C$ 的同伦型不变量。这里所定义的是定界余链的
规范等价\liftIndex{I152}\liftIndex{I153}。

\subsection{基本性质与范畴 $\mathfrak{HA}_\infty$}
\label{subsec:4-3-1}

令 $C$ 为滤过 $A_\infty$ 代数，令 $\mathfrak C$ 为
$[0,1]\times C$ 的一个模型。

\begin{definition}{4.3.1}\label{def:4-3-1}
令 $b_0,b_1\in\widehat{\mathcal M}(C)$。若存在
$\widetilde b\in\widehat{\mathcal M}(\mathfrak C)$，使得
\[
 \mathrm{Eval}_{s=0*}(\widetilde b)=b_0,
 \qquad
 \mathrm{Eval}_{s=1*}(\widetilde b)=b_1,
\]
就称 $b_0$ 与 $b_1$ \emph{$\mathfrak C$-规范等价}，并记作
$b_0\sim_{\mathfrak C}b_1$。称这样的 $\widetilde b$ 为从 $b_0$ 到
$b_1$ 的一个同伦。
\end{definition}

\begin{remark}{4.3.2}\label{rem:4-3-2}
集合 $\widehat{\mathcal M}(C)$ 及其上的等价关系 $\sim$ 只对滤过
$A_\infty$ 代数定义。事实上，在无滤过 $A_\infty$ 代数中，由于
$\widehat d(e^b)$ 一般是无穷和，Maurer--Cartan 方程
$\widehat d(e^b)=0$ 没有意义，因而所有定界余链组成的集合
$\widehat{\mathcal M}(C)$ 也没有意义。不过，仍可把
$\widehat{\mathcal M}(C)$ 定义为 Artin 环范畴上的一个函子
（\textbf{[Schl68]}）；例如见 \textbf{[Fuk03II]} 第 2.1 节。
\end{remark}

\begin{lemma}{4.3.3}\label{lem:4-3-3}
关系 $\sim_{\mathfrak C}$ 不依赖于 $[0,1]\times C$ 的模型
$\mathfrak C$ 的选择。
\end{lemma}

\begin{proof}
证明与引理 4.2.36 的证明相同。
\end{proof}

从现在起，以 $\sim$ 代替 $\sim_{\mathfrak C}$，并以“规范等价”代替
“$\mathfrak C$-规范等价”。

\begin{lemma}{4.3.4}\label{lem:4-3-4}
定界余链集合上的规范等价关系 $\sim$ 是等价关系。
\end{lemma}

\begin{proof}
证明与命题 4.2.37 的证明相同。
\end{proof}

下面更显式地描述规范等价。先取
$\mathfrak C=\operatorname{Poly}([0,1],C)$。

\begin{proposition}{4.3.5}\label{prop:4-3-5}
$b_0\in\widehat{\mathcal M}(C)$ 与
$b_1\in\widehat{\mathcal M}(C)$ 规范等价，当且仅当存在
\[
 \widetilde b=(b(s),c(s))
 \in\operatorname{Poly}([0,1],C)^1,
\]
使得
\begin{align}
 \frac{\partial b(s)}{\partial s}
 &=\delta_{b(s)}c(s),
 \tag{4.3.6.1}\label{eq:4-3-6-1}\\
 b(0)&=b_0,\qquad b(1)=b_1.
 \tag{4.3.6.2}\label{eq:4-3-6-2}
\end{align}
\end{proposition}

\begin{proof}
令 $\widetilde b$ 为从 $b_0$ 到 $b_1$ 的同伦，并写成
$\widetilde b=(b(s),c(s))$。条件
$\widetilde b\in\widehat{\mathcal M}
(\operatorname{Poly}([0,1],C))$ 等价于 \eqref{eq:4-3-6-1}，而条件
$\mathrm{Eval}_{s=0*}(\widetilde b)=b_0$、
$\mathrm{Eval}_{s=1*}(\widetilde b)=b_1$ 等价于
\eqref{eq:4-3-6-2}。
\end{proof}

在微分分次代数 $(C;d,\circ)$ 的情形，若 $d$ 是导子、$\circ$ 是结合积，
则可证明：若 $b$ 与 $b'$ 规范等价，就存在 $g$，使得
\[
 b'=g^{-1}\circ d(g)+g^{-1}\circ b\circ g.
\]
换言之，$b'$ 按通常意义与 $b$ 规范等价。由于后文不需要这个结果，这里不予
证明（见 \textbf{[Fuk03II]} 引理 2.2.2）。还可指出下面的事实。

\begin{lemma}{4.3.7}\label{lem:4-3-7}
若 $(b(s),c(s))$ 满足 \eqref{eq:4-3-6-1}，并且 $b(0)$ 是定界余链，
则对每个 $s$，$b(s)$ 也是定界余链。
\end{lemma}

\begin{proof}
利用 \eqref{eq:4-3-6-1} 计算得
\[
 \frac{\partial}{\partial s}\mathfrak m(e^{b(s)})
 =\mathfrak m\bigl(e^{b(s)},\delta_{b(s)}(c(s)),e^{b(s)}\bigr)
 =\delta_{b(s)}\bigl(\delta_{b(s)}(c(s))\bigr)=0.
\]
引理 4.3.7 随即成立。
\end{proof}

\begin{remark}{4.3.8}\label{rem:4-3-8}
命题 4.3.5 的结论类似于微分分次代数文献中的相应定义（见
\textbf{[GoMi88, Sche98, Kon03]}）。\textbf{[Kon03]} 对
$L_\infty$ 代数提出了类似定义。不过，作者没有在文献中找到引理 4.3.3、
4.3.4 的证明。
\end{remark}

下面取 $\mathfrak C=C^{[0,1]}$。

\begin{proposition}{4.3.9}\label{prop:4-3-9}
两个定界余链 $b_0,b_1$ 彼此规范等价，当且仅当存在
$c\in C[1]^{-1}$，使得
\begin{equation}
 b_1-b_0=\mathfrak m(e^{b_0}ce^{b_1}).
\tag{4.3.10}\label{eq:4-3-10}
\end{equation}
\end{proposition}

\begin{proof}
若 $\widetilde b\in\widehat{\mathcal M}(C^{[0,1]})$ 是从 $b_0$ 到
$b_1$ 的同伦，则可写成
$\widetilde b=(b_0,c,b_1)$。方程 \eqref{eq:4-3-10} 等价于条件
$\widetilde b\in\widehat{\mathcal M}(C^{[0,1]})$。
\end{proof}

\begin{definition}{4.3.11}\label{def:4-3-11}
令 $(C,\mathfrak m)$ 为滤过 $A_\infty$ 代数。定义
$\mathcal M(C,\mathfrak m)$ 为
$\widehat{\mathcal M}(C,\mathfrak m)$ 的所有 $\sim$-等价类组成的集合。
为简便起见，有时分别把它们写作
$\mathcal M(C)$、$\widehat{\mathcal M}(C)$。
\end{definition}

\begin{definition}{4.3.12}\label{def:4-3-12}
对辛流形 $M$ 的相对 spin 拉格朗日子流形 $L$，第 3.5 节配置了滤过
$A_\infty$ 代数
$(C(L;\Lambda_{0,\mathrm{nov}}),\mathfrak m)$。记
\[
 \mathcal M(L)=
 \mathcal M(C(L;\Lambda_{0,\mathrm{nov}}),\mathfrak m).
\]
\end{definition}

本章的主要结果之一是，
$\mathcal M(C(L;\Lambda_{0,\mathrm{nov}}),\mathfrak m)$ 不依赖于构造中
涉及的各种选择（定理 4.1.3）。下面的定理是其证明的代数部分。令
\[
 \mathfrak{HA}_\infty=
 \bigl[\text{滤过 $A_\infty$ 代数；滤过 $A_\infty$ 同态的同伦类}\bigr]
\]
表示这样一个范畴：对象是滤过 $A_\infty$ 代数，态射是滤过
$A_\infty$ 同态的同伦类。

\begin{theorem}{4.3.13}\label{thm:4-3-13}
对应
$(C,\mathfrak m)\longmapsto\mathcal M(C,\mathfrak m)$ 定义了协变函子
\[
 \mathfrak{HA}_\infty\longrightarrow[\mathrm{set}].
\]
\end{theorem}

下面的推论 4.3.14 是 W. Goldman 与 J. Millson
\textbf{[GoMi88,90]} 的一个结果在滤过 $A_\infty$ 代数中的类似物。

\begin{corollary}{4.3.14}\label{cor:4-3-14}
若
$\mathfrak f\colon(C_1,\mathfrak m^1)\to(C_2,\mathfrak m^2)$
是同伦等价，则它诱导双射
\[
 \mathfrak f_*\colon
 \mathcal M(C_1,\mathfrak m^1)\longrightarrow
 \mathcal M(C_2,\mathfrak m^2).
\]
\end{corollary}

为证明定理 4.3.13，只须证明下面的引理 4.3.15。回忆，对滤过
$A_\infty$ 同态
$\mathfrak f\colon(C_1,\mathfrak m^1)\to(C_2,\mathfrak m^2)$，引理
3.6.36 已经通过
\[
 \mathfrak f_*(b)=\mathfrak f(e^b)
\]
定义了映射
$\mathfrak f_*\colon\widehat{\mathcal M}(C_1,\mathfrak m^1)
\to\widehat{\mathcal M}(C_2,\mathfrak m^2)$。

\pagebreak[3]
\begin{lemma}{4.3.15}\label{lem:4-3-15}
\begin{enumerate}[label=\textup{(\arabic*)},leftmargin=2.8em]
 \item 若 $b_0\sim b_1$ 于
 $\widehat{\mathcal M}(C_1,\mathfrak m^1)$，则
 $\mathfrak f_*(b_0)\sim\mathfrak f_*(b_1)$。
 \item 若 $\mathfrak f$ 同伦于 $\mathfrak f'$，则对任意
 $b\in\widehat{\mathcal M}(C_1,\mathfrak m^1)$，都有
 $\mathfrak f_*(b)\sim\mathfrak f'_*(b)$。
\end{enumerate}
\end{lemma}

\begin{proof}
令 $\mathfrak C_i$（$i=1,2$）为 $[0,1]\times C_i$ 的一个模型。

(1) 把定理 4.2.34 用于 $\mathfrak f$，得到滤过 $A_\infty$ 同态
$\mathfrak F\colon\mathfrak C_1\to\mathfrak C_2$，满足
\[
 \mathrm{Eval}_{s=s_0}\circ\mathfrak F
 =\mathfrak f\circ\mathrm{Eval}_{s=s_0}.
\]
令 $\widetilde b\in\widehat{\mathcal M}(\mathfrak C_1)$ 为从 $b_0$
到 $b_1$ 的同伦。于是 $\mathfrak F_*(\widetilde b)$ 是从
$\mathfrak f_*(b_0)$ 到 $\mathfrak f_*(b_1)$ 的同伦。因此
$\mathfrak f_*(b_0)\sim\mathfrak f_*(b_1)$。

(2) 令 $\mathfrak H\colon C_1\to\mathfrak C_2$ 为从 $\mathfrak f$
到 $\mathfrak f'$ 的同伦。则 $\mathfrak H_*(b)$ 是定界余链，并满足
\[
 \mathrm{Eval}_{s=0*}\bigl(\mathfrak H_*(b)\bigr)=\mathfrak f_*(b),
 \qquad
 \mathrm{Eval}_{s=1*}\bigl(\mathfrak H_*(b)\bigr)=\mathfrak f'_*(b).
\]
故 $\mathfrak f_*(b)\sim\mathfrak f'_*(b)$。
\end{proof}

在下一个命题的证明中，将用到定理 4.2.45；该定理要到第 4.5 节才会证明。

\begin{proposition}{4.3.16}\label{prop:4-3-16}
若 $b_0\sim b_1$ 于 $\widehat{\mathcal M}(C,\mathfrak m)$，则存在严格
同伦等价
\[
 \mathfrak f\colon(C,\mathfrak m^{b_0})
 \longrightarrow(C,\mathfrak m^{b_1}),
\]
满足
$\overline{\mathfrak f}\equiv\mathrm{id}\pmod{\Lambda^+_{0,\mathrm{nov}}}$。
这里，$\mathfrak m^{b_i}$ 是定义 3.6.9 所定义的形变
$A_\infty$ 结构。
\end{proposition}

\begin{proof}
令 $\mathfrak C$ 为 $[0,1]\times C$ 的一个模型，并令 $\widetilde b$
如定义 4.3.1。于是 $\mathrm{Eval}_{s=0}$ 与
$\mathrm{Eval}_{s=1}$ 定义严格滤过 $A_\infty$ 同态
\[
 \mathrm{Eval}_{s=0}\colon
 (\mathfrak C,\mathfrak m^{\widetilde b})
 \longrightarrow(C,\mathfrak m^{b_0}),
 \qquad
 \mathrm{Eval}_{s=1}\colon
 (\mathfrak C,\mathfrak m^{\widetilde b})
 \longrightarrow(C,\mathfrak m^{b_1}).
\]
利用定理 4.2.45，可以证明二者都是严格同伦等价，由此得到
$\mathfrak f$。将在第 4.5 节末检验
$\overline{\mathfrak f}\equiv\mathrm{id}\pmod{\Lambda^+_{0,\mathrm{nov}}}$。
\end{proof}

\begin{remark}{4.3.17}\label{rem:4-3-17}
对任意 $b\in\widehat{\mathcal M}(C,\mathfrak m)$，滤过
$A_\infty$ 代数 $(C,\mathfrak m^b)$ 都同伦等价于
$(C,\mathfrak m)$。事实上，定义
\[
 \mathfrak i^b\colon(C,\mathfrak m^b)\longrightarrow(C,\mathfrak m)
\]
为
\[
 \mathfrak i^b_0(1)=b,\qquad
 \mathfrak i^b_1=\mathrm{id},\qquad
 \mathfrak i^b_k=0\quad(k\geq2).
\]
由定理 4.2.45，它是同伦等价。不过要注意，$\mathfrak i^b$ 在定义
3.2.29 的意义下并不严格。
\end{remark}

下面的命题是推论 5.2.40 在 $L=L^{(0)}=L^{(1)}$ 时的特例。

\begin{proposition}{4.3.18}\label{prop:4-3-18}
令
$b_0,b_1,b'_0,b'_1\in\widehat{\mathcal M}(C,\mathfrak m)$，并假定
$b_0\sim b_1$ 且 $b'_0\sim b'_1$。则存在典范同构
\[
 H^*(C;\delta_{b_0,b'_0})
 \cong H^*(C;\delta_{b_1,b'_1}).
\]
\end{proposition}

\begin{proof}
取定界余链 $\widetilde b,\widetilde b'\in C^{[0,1]}$，使得对
$i=0,1$ 都有
\[
 \mathrm{Eval}_{s=i}(\widetilde b)=b_i,
 \qquad
 \mathrm{Eval}_{s=i}(\widetilde b')=b'_i.
\]
考虑下面的分次模正合列：
\[
 0\longrightarrow C[-1]\oplus C
 \longrightarrow C^{[0,1]}
 \longrightarrow C^{[0,1]}/(C[-1]\oplus C)
 \longrightarrow0,
\]
其中第二个同态是包含映射。注意，$C[-1]\oplus C$ 对
$\delta_{\widetilde b,\widetilde b'}$ 与 $\mathfrak m_{1,0}$ 都封闭；
$\overline C[-1]\oplus\overline C$ 是
$(\overline C^{[0,1]},\overline{\mathfrak m}_1)$ 的子复形。还容易看出
$(\overline C[-1]\oplus\overline C,\overline{\mathfrak m}_1)$ 是无环
复形，因而 $(C[-1]\oplus C,\mathfrak m_{1,0})$ 也是无环的。下面证明
$(C[-1]\oplus C,\delta_{\widetilde b,\widetilde b'})$ 无环。

令
\[
 c=\sum_{j=0}^{\infty}c_j e^{\nu_j}T^{\lambda_j}
 \in C[-1]\oplus C
\]
为 $\delta_{\widetilde b,\widetilde b'}$-闭上链。为简单起见，假设
$\lambda_j$ 严格单调递增：
$0=\lambda_0<\lambda_1<\lambda_2<\cdots$。由于
\[
 \delta_{\widetilde b,\widetilde b'}
 \equiv\overline{\mathfrak m}_1
 \pmod{\Lambda^+_{0,\mathrm{nov}}},
\]
$c_0\in\overline C[-1]\oplus\overline C$ 是
$\overline{\mathfrak m}_1$-闭上链。回忆
$\overline C[-1]\oplus\overline C$ 关于
$\overline{\mathfrak m}_1$ 无环，故存在
$a_0\in\overline C[-1]\oplus\overline C$，使得
$c_0=\overline{\mathfrak m}_1a_0$。置
\[
 c'=c-\delta_{\widetilde b,\widetilde b'}(a_0),
\]
则 $c'$ 仍是 $\delta_{\widetilde b,\widetilde b'}$-闭上链。（回忆按本书
对 $\Lambda_{0,\mathrm{nov}}$ 的假设，$\mu_0=0$。）$c'$ 的能级严格高于
$c$；也就是说，对某个 $\lambda>0$，有
$c'\in F^\lambda\mathfrak C$。利用有隙条件，可以反复施行上述步骤，把
$c$ 写成一个 $\delta_{\widetilde b,\widetilde b'}$-上边界。因此，
$C[-1]\oplus C$ 关于 $\delta_{\widetilde b,\widetilde b'}$ 无环。

商复形 $C^{[0,1]}/(C[-1]\oplus C)$ 连同由
$\delta_{\widetilde b,\widetilde b'}$ 诱导的上边界算子，自然同构于复形
$(C,\delta_{b_0,b'_0})$，所以
\[
 H^*(C^{[0,1]};\delta_{\widetilde b,\widetilde b'})
 \cong H^*(C;\delta_{b_0,b'_0}).
\]
同理，
\[
 H^*(C^{[0,1]};\delta_{\widetilde b,\widetilde b'})
 \cong H^*(C;\delta_{b_1,b'_1}).
\]
这就证明了命题 4.3.18。
\end{proof}

特别地，命题 4.3.18 蕴含：弗洛尔上同调群
\[
 HF((L,b),(L,b');\Lambda_{0,\mathrm{nov}})
\]
只依赖于 $b,b'$ 的规范等价类\liftIndex{I154}。第 5.3、5.4 节将进一步
研究 $\mathcal M(C,\mathfrak m)$。
